\documentclass{article}
\usepackage{amsmath,amssymb}
\usepackage[margin=1in]{geometry}
\allowdisplaybreaks
\usepackage[UTF8]{ctex}
\begin{document}

\section*{最优的返利合同设计，结合内生预测技术}

\subsection*{无预测合同：目标与约束条件}
制造商选择了一个单一的合同$(q_N,r_N,t_N)$:
\begin{equation}
\max_{q_N,r_N,t_N}\left\{t_N-cq_N-r_NE\min(q_N,D_N)\right\}.
\tag{OBJ--N}
\end{equation}
\begin{align}
\text{s.t.}\quad R^r(N,N)&\geq \lambda\max(R^r(H,N),0)
 +(1-\lambda)\max(R^r(L,N),0)-k,\tag{IC}\\
R^r(N,N)&\geq 0.\tag{IR}
\end{align}
这些是第5.1节中的目标和主体部分的内容（打印在~1658上）。第一个约束条件是不允许进行预测；第二个约束条件则是确保无需进行预测即可接受该方案。

\subsection*{预测菜单：目标与约束条件}
制造商可以选择“两合同”菜单$\{(q_H,r_H,t_H),(q_L,r_L,t_L)\}$，其中合同$S\in\{H,L\}$适用于那些能够预测并观察信号$S$的零售商。
\begin{equation}
\begin{split}
\max_{(q_H,r_H,t_H),(q_L,r_L,t_L)}\bigl\{
&\lambda[t_H-cq_H-r_HE\min(q_H,D_H)]\\
&+(1-\lambda)[t_L-cq_L-r_LE\min(q_L,D_L)]\bigr\}.
\end{split}
\tag{OBJ--F}
\end{equation}
\begin{align}
\text{s.t.}\quad R^r(H,H)&\geq R^r(H,L),\tag{IC1}\\
R^r(L,L)&\geq R^r(L,H),\tag{IC2}\\
R^r(H,H)&\geq 0,\tag{IR1}\\
R^r(L,L)&\geq 0,\tag{IR2}\\
\lambda R^r(H,H)+(1-\lambda)R^r(L,L)-k
&\geq R^r(N,H),\tag{IC3}\\
\lambda R^r(H,H)+(1-\lambda)R^r(L,L)-k
&\geq R^r(N,L),\tag{IC4}\\
\lambda R^r(H,H)+(1-\lambda)R^r(L,L)-k
&\geq 0.\tag{IR3}
\end{align}
这是第5.2节的内容（打印页码为~1659–1660）。约束条件（IC1）到（IR2）用于在观察到信号后选择参与方；之后的成本预测则被忽略。约束条件（IC3）到（IR3）则用于之前的预测和参与决策。这两个部分属于不同的方案，并非同时对单一菜单施加的约束条件。下面将比较这两种方案的最优利润。

\subsection*{继承下来的定义、领域以及信息结构保持不变。}
固定输入包括零售价格$p$、单位生产成本$c$、预测成本$k$、信号概率$\lambda$，以及需求分布。生产发生在销售季节之前；未售出的库存几乎没有回收价值。零售商的保留利润被设定为零。需求是外生的，模型仅涉及一种产品和一个时期，并且采用二元预测决策方式。

无条件要求$D_N$是一个非负随机变量，其分布为$F_N$。预测成本$k$，并私下透露一个信号：
\[
\Pr(S=H)=\lambda,\qquad \Pr(S=L)=1-\lambda.
\]
条件需求$D_S$是非负的，并且其分布为$F_S$，其中$S\in\{H,L\}$为某个特定值。定义分布关系如下：
\[
F_N(x)=\lambda F_H(x)+(1-\lambda)F_L(x),
\]
\[
F_H(x)\leq F_L(x)\qquad\text{for all }x\geq 0,
\]
\[
\bar F_S(x)\equiv 1-F_S(x),\qquad S\in\{H,L,N\}.
\]
本文假设$F_S(\cdot)$的分布是严格递增的，以便于表述。同时指出，即使这一假设被放宽，其结论仍然成立。条件分布和预测成本都是众所周知的知识；只有零售商能够知道她是否进行了预测，如果进行了预测，那么她所得到的信号就是实际值。这里的$E$表示在由随机变量所表示的需求分布下的期望值。

制造商的返利转账功能具有已打印的定义域和陪域。
\[
T^r(q,r):\mathbb R^+\times\mathbb R^+\longrightarrow\mathbb R^+.
\]
数量$q$是零售商的下单数量，$r$是每售出一个单位商品时零售商给予最终消费者的返利，而$T^r(q,r)$则是零售商支付给制造商的转让费用。因此，合同中的各项数值都具有非负的取值范围。
\[
q_C\geq 0,\qquad r_C\geq 0,\qquad t_C\geq 0,
\]
在“无预测”模块中保留$C=N$，在“预测”模块中保留$C\in\{H,L\}$。特别是，$t_C$是一个总转移量，而不是单位价格。制造商自行选择这些坐标；零售商则可以选择是否进行预测，以及接受哪个报价或拒绝该报价。

第3节（打印在~1657处）给出了双合同菜单的表示方式。
\[
T^r(q,r)=
\begin{cases}
t_L,&q\leq q_L\text{ and }r\leq r_L,\\
t_H,&q_L<q\leq q_H\text{ and }r\leq r_H,\\
\infty,&\text{otherwise},
\end{cases}
\qquad (q,r)\in\mathbb R^+\times\mathbb R^+,
\]
并将其紧凑地写成：
\[
\{(q_H,r_H,t_H),(q_L,r_L,t_L)\}.
\]
第5.2节中的“揭示原则”构建方式将注意力限制在两种合同上，因为存在两种可能的信号。当制造商希望无需预测即可接受订单时，就会使用单一合同。论文中提到，转让价格通常会随着数量的增加而上升，同时还有折扣政策也会更加优惠，这些描述性内容并不构成任何额外的约束条件，也不属于打印出来的优化块中的内容。

事件的顺序如下：制造商提供菜单；零售商决定是否进行预测，如果进行预测，则私下观察信号；零售商选择合同或拒绝报价；一旦接受，制造商就会生产所订购的数量，并收到转账；需求得到满足，然后支付回扣。

零售商的预期利润，不包括任何预测成本，为
\[
R^r(S,C)=(p+r_C)E\min(q_C,D_S)-t_C,
\qquad S\in\{H,L,N\},
\]
其中，$C$表示适用制度下的各种合同。上标$r$表示返利合同形式；$S$表示需求信息，而$C$则表示所选定的合同。

在单一报价下，先进行预测，然后做出最优的接受决策。
\[
\lambda\max(R^r(H,N),0)
 +(1-\lambda)\max(R^r(L,N),0)-k.
\]
无条件利润恒等式是
\[
\lambda R^r(H,N)+(1-\lambda)R^r(L,N)=R^r(N,N).
\]
如第5.1节所述，只有当低信号使得接受决策变得不具吸引力时，预测才具有价值。
\[
R^r(L,N)<0.
\]
这最后一个关系描述了那种情况；它并不是对每个不预测合同都施加的额外约束条件。

\subsection*{由“optimal-contract”特征所定义的术语}
在命题2和命题3中使用的综合系统量，其定义见于第4节。
\[
q_N^I=\bar F_N^{-1}(c/p),\qquad
q_S^I=\bar F_S^{-1}(c/p),\quad S\in\{H,L\}.
\]
逆符号按原样保留，不添加任何通用的逆符号约定。第5.1节对此进行了定义。
\[
\Gamma(q)\equiv(1-\lambda)p\int_0^q
 [\bar F_N(x)-\bar F_L(x)]\,dx.
\]
它指出，$q_N^I\geq q_L^I$，并且对于任何$q\in(q_L^I,q_N^I)$来说，
\[
\bar F_N(q)>c/p,\qquad \bar F_L(q)<c/p,
\]
\[
\Gamma'(q)=(1-\lambda)p[\bar F_N(q)-\bar F_L(q)]>0.
\]
因此，在返利描述中使用的逆运算$\Gamma^{-1}(\cdot)$被定义为在…上良好定义的。
\[
[\Gamma(q_L^I),\Gamma(q_N^I)].
\]
后来引入的不同的逆向约定不适用于退货合同。

第5.2节中定义了
\[
\Delta(q)\equiv E[\min(q,D_H)-\min(q,D_L)],
\]
在收到有利信号后，预计销售量会有所增加。

\subsection*{最优的无需预测的合同及其利润}
命题2给出了一个最优的契约。
\[
(q_N^*,r_N^*,t_N^*)=
\begin{cases}
\left(q_L^I,0,pE\min(q_L^I,D_N)
 -\dfrac{\Gamma(q_L^I)-k}{1-\lambda}\right),
 &k\leq\Gamma(q_L^I),\\[6pt]
\left(\Gamma^{-1}(k),0,pE\min(\Gamma^{-1}(k),D_N)\right),
 &k\in(\Gamma(q_L^I),\Gamma(q_N^I)),\\[6pt]
\left(q_N^I,0,pE\min(q_N^I,D_N)\right),
 &k\geq\Gamma(q_N^I).
\end{cases}
\]
特别是，$r_N^*=0$和$q_N^*\leq q_N^I$。制造商所预测的最佳预期利润为：
\begin{equation}
\mathcal M_N^r=
\begin{cases}
pE\min(q_L^I,D_N)-cq_L^I
 -\dfrac{\Gamma(q_L^I)-k}{1-\lambda},
 &k\leq\Gamma(q_L^I),\\[6pt]
pE\min(\Gamma^{-1}(k),D_N)-c\Gamma^{-1}(k),
 &k\in(\Gamma(q_L^I),\Gamma(q_N^I)),\\[6pt]
pE\min(q_N^I,D_N)-cq_N^I,
 &k\geq\Gamma(q_N^I).
\end{cases}
\tag{1}
\end{equation}
在最佳状态下，文本还指出：
\[
R^r(N,N)=\frac{\Gamma(q_L^I)-k}{1-\lambda}>0
\qquad\text{when }k<\Gamma(q_L^I).
\]

\subsection*{最优预测合约及其利润}
如果预测问题（OBJ--F)--(IR3）没有解决方案，那么论文建议制造商不要进行预测。第3个命题则针对该问题有解决方案的情况进行了讨论。
\begin{align*}
q_L^*&=\operatorname*{arg\,max}_{q_L\geq0}
 \left\{p\int_0^{q_L}[\bar F_L(x)-\lambda\bar F_H(x)]\,dx
 -(1-\lambda)cq_L\right\},\\
r_L^*&=0,\\
t_L^*&=pE\min(q_L^*,D_L),\\
q_H^*&=q_H^I,\\
r_H^*&=k/\lambda(1-\lambda)\Delta(q_H^I),\\
t_H^*&=(p+r_H^*)E\min(q_H^*,D_H)
 -p\Delta(q_L^*)-k/\lambda.
\end{align*}
\noindent\emph{印刷式注释。}上述$r_H^*$的折扣行保留了在命题3中印刷的未加括号的实线表达式（物理PDF第~7页）。来源并未对分母进行分隔；没有添加额外的括号或修正分数。

这些是解决方案的描述，并非对（OBJ--F)--(IR3）中任意菜单所施加的额外约束。随附的文本中说明了这些内容。
\[
q_L^*\leq q_L^I,\qquad q_H^*=q_H^I,
\]
在所描述的菜单中，零售商在扣除预测成本之前的利润是
\[
R^r(L,L)=0,\qquad R^r(H,H)=p\Delta(q_L^*)+k/\lambda.
\]
其预期的收益满足所关联的方程组。
\begin{align}
\lambda R^r(H,H)+(1-\lambda)R^r(L,L)-k
 &=\lambda R^r(H,L)\tag{2}\\
 &=R^r(N,L)\tag{3}\\
 &=\lambda p\Delta(q_L^*).\tag{4}
\end{align}
论文指出，当$F_L(0)=0$和$F_H(x)<F_L(x)$满足某些足够小的$x$时，这种利润是严格正的。这是得出“严格正利润”结论的充分条件，而不是对优化模型的其他限制条件。

在预测模式下，制造商的最佳预期利润是
\begin{equation}
\begin{split}
\mathcal M_F^r={}&\lambda\left[
 p\int_0^{q_H^*}\bar F_H(x)\,dx-cq_H^*\right]\\
&+(1-\lambda)\left[
 p\int_0^{q_L^*}\frac{\bar F_L(x)-\lambda\bar F_H(x)}{1-\lambda}\,dx
 -cq_L^*\right]-k.
\end{split}
\tag{5}
\end{equation}

\subsection*{内生性预测机制的选择}
5.3节以及印刷版中的~1662开头部分，比较了两种最优利润情况：
\[
\mathcal M^r=\max(\mathcal M_F^r,\mathcal M_N^r).
\]
在以下条件下，预测问题才有解决方案：制造商在以下情况下才会进行预测：当且仅当……
\[
\mathcal M_F^r>\mathcal M_N^r.
\]
否则，制造商会提供“无预测”合同。具体来说，该规则采用严格的比较方式，将“相等”视为“无预测”。论文指出，$\mathcal M_F^r$会随着$k$而严格减少，而$\mathcal M_N^r$则会随着$k$而增加。
\[
\left.\mathcal M_F^r\right|_{k=0}
 \geq \left.\mathcal M_N^r\right|_{k=0}.
\]
因此，它确定了一个满足条件的阈值$k^r\geq0$。
\[
\mathcal M_F^r>\mathcal M_N^r
\quad\Longleftrightarrow\quad k<k^r.
\]
对于$k<k^r$，制造商提供了“命题3”选项；否则，制造商则提供“命题2”单一合同选项。在本文中，并未提供$k^r$的明确公式。

\paragraph{来源范围与资格要求。}
操作来源是 Taylor 和 Xiao 的《\emph{零售商预测激励措施：返利与退货}》，以及《\emph{管理科学}》第 55(10) 期，2009 年，第 3-5 节，印刷页码为~1656-1662页。其中，(OBJ--N) 和 (OBJ--F) 标签分别指代源代码中标记为 (OBJ) 的两个目标。退货合同、线性转移限制以及后续的稳健性扩展内容并不属于本模型的范畴。

打印出的地图$T^r:\mathbb R^+\times\mathbb R^+\to\mathbb R^+$及其后续分配$T^r=\infty$被保留下来，即使它们的 Codomain 并不匹配。所提供的原始模型段落中并未明确提及严格的界限$0<\lambda<1$或$0<c<p$，也没有关于消失的分母或逆边界情况的任何规定；这些内容都没有被添加进来。论文中提到了一个电子版的补充材料，但该材料并未提供，也未被验证；这里也没有提及任何仅适用于该材料的条件或替代方案。

\end{document}